\documentclass[10pt,a4paper]{amsart}
\usepackage[margin=19mm]{geometry}
\usepackage{amsmath,amssymb,mathtools}
\usepackage[scr=rsfs,cal=euler]{mathalfa}
\usepackage{xfrac}
\usepackage[unicode,hidelinks,bookmarks=false]{hyperref}
\newcommand{\ie}{i.e.\ }
\newcommand{\cf}{cf.\ }
\newcommand{\resp}{respectively\ }
\newcommand{\Subsection}{\subsection}
\providecommand{\nobreakdash}{\nobreak\hbox{-}}
\setlength{\emergencystretch}{2em}
\multlinegap0pt
\usepackage{amssymb}
\usepackage{tikz}
\usepackage{dsfont}
\newtheorem{lem}{Lemma}
\newtheorem{df}{Definition}
\title{Universal Approximation of Operators with Transformers and Neural Integral Operators}
\author{Emanuele Zappala$^*$, Maryam Bagherian}
\date{Source: arXiv:2409.00841v3 (CC BY 4.0)}
\begin{document}
\maketitle

\noindent\emph{Source and licence.} Universal Approximation of Operators with Transformers and Neural Integral Operators. Version arXiv:2409.00841v3 (CC BY 4.0), distributed by arXiv under the Creative Commons Attribution 4.0 International licence, http://creativecommons.org/licenses/by/4.0/, as recorded in the arXiv licence metadata for this exact version and shown on the abstract page https://arxiv.org/abs/2409.00841v3. Full licence notice: \texttt{licenses/arxiv-2409.00841-CC-BY-4.0.txt}.\par\smallskip

\noindent\textit{Editorial scope.} Counted unit: the article's lem lem:equiv\_gen (source0.tex lines 403--405). The article's complete original proof of it is delivered with it (source0.tex lines 406--441, one \texttt{proof} environment of 36 lines). No mathematics is written, repaired or re-derived: the statement and the proof are the article's own lines, in the article's own notation, and no retained line is substituted or re-flowed.  Retained uncounted as the source-local context the proof needs: lines 351--357; lines 382--401.  Nothing was omitted from any retained span, so the package carries no omission record.  No retained line contains a citation, so the wrapper carries no bibliography and no bibliography entry.  No retained line contains a figure environment, an image-inclusion command or drawing code, so no image is delivered and the package declares no asset dependency.  Editorial formatting only, in addition: an AMS \texttt{amsart} preamble that restates the article's own declarations and macros used by the retained lines, namely the environments \texttt{lem}, \texttt{df} on separate counters, together with the standard packages those lines need.  The retained environments print in this document as Lemma 1, Df 1 in order of appearance, whereas the article prints the same items with the numbering of its own full document.  The complete native source of the article as distributed in the arXiv e-print for this exact version is bundled unchanged as \texttt{sources/164586/source0.tex}.
\begin{lem}\label{lem:int_scheme}
        Let $T$ and $X$ be as above. Then, we can find an integration scheme, $\mathcal T$ such that 
        \begin{eqnarray}
            \|T(y) - \mathcal T(y)\|_{\infty} < \epsilon,
        \end{eqnarray}    
        for all $y\in \bar {\mathbb B}_\rho$, where $\|\cdot \|_{\infty}$ denotes the uniform norm. 
\end{lem}

\begin{df}\label{def:actions}
    {\rm 
Let $\sigma$ denote a permutation of $n$ elements, i.e. $\sigma\in \Sigma_n$ where $\Sigma_n$ is the permutation group on $n$ elements. Then, $\sigma$ acts on the tuples of type $(z_1, \ldots, z_n, x_1, \ldots , x_n)\in \mathbb R^{2n}$ through the induced diagonal permutation representation, which just reorders the entries $z_1, \ldots , z_n$ and $x_1, \ldots , x_n$ according to the permutation $\sigma$: 
            $$(z_1, \ldots, z_n, x_1, \ldots, x_n)\cdot_1 \sigma := (z_{\sigma(1)}, \ldots, z_{\sigma(n)}, x_{\sigma(1)}, \ldots, x_{\sigma(n)}).$$ 
            We indicate by $\cdot_2$ the permutation action of $\Sigma_n$ on $n$-tuples as
            $$
                    \begin{bmatrix}
                                        x_1\\
                                        \vdots\\
                                        x_n
                                 \end{bmatrix}\cdot_2  \sigma = 
                                 \begin{bmatrix}
                                       x_{\sigma(1)}\\
                                       \vdots\\
                                       x_{\sigma(n)}
                                 \end{bmatrix}.
            $$
            We will denote by $\cdot_2$ also the permutation action on tuples such as $(x_1,\ldots, x_n)$ to shorten notation. 
    }
\end{df}

\begin{lem}\label{lem:equiv_gen}
    The integration scheme from Lemma~\ref{lem:int_scheme} induces a continuous function $F: \mathbb R^n\times [a,b]^n \longrightarrow \mathbb R^n$. Moreover, $F$ is equivariant with respect to the actions defined in Definition~\ref{def:actions}.
\end{lem}

\begin{proof}
        We set $F^i(z_1, \ldots, z_n,x_1,\ldots, x_n) := \sum_{k=1}^n G(z_k,x_i,t_k)w_k$, for each $i=1,\ldots , n$, where the points $t_1, \ldots, t_n$ are determined by the integration scheme defined in Lemma~\ref{lem:int_scheme}. 
        We can therefore define $F$ as the column vector $[F^i]_{i=1}^n$, as a function of the tuple $(z_1, \ldots, z_n, x_1, \ldots, x_n)$. Since each function $F^i$ is continuous with respect to its entries $z_k$ and $x_i$, it follows that the function $F$ is also continuous. 
        
        To show that the function $F$ just defined is permutation equivariant, using the same notations introduced above, we compute
        \begin{eqnarray*}
                F( (z_1, \ldots, z_n, x_1, \ldots, x_n)\cdot_1 \sigma) &=& F(z_{\sigma(1)}, \ldots, z_{\sigma(n)}, x_{\sigma(1)}, \ldots, x_{\sigma(n)})\\
                &=& \begin{bmatrix}
             F^{\sigma(1)}(z_{\sigma(1)}, \ldots, z_{\sigma(n)},x_{1}, \ldots, x_{n})\\
                \vdots\\
            F^{\sigma(n)}(z_{\sigma(1)}, \ldots, z_{\sigma(n)},x_{1}, \ldots, x_{n})
              \end{bmatrix}\\
          &=&  \begin{bmatrix}
                    \sum_{k=1}^n G(z_{\sigma(k)},x_{\sigma(1)},t_{\sigma(k)})w_{\sigma(k)}\\
                    \vdots\\
                    \sum_{k=1}^n G(z_{\sigma(k)},x_{\sigma(n)},t_{\sigma(k)})w_{\sigma(k)}
               \end{bmatrix}\\
          % &=&  \begin{bmatrix}
          %           \sum_{k=1}^n G(z_{\sigma(k)},t_{\sigma(n+1)})w_{\sigma(k)}\\
          %           \vdots\\
          %           \sum_{k=1}^n G(z_{\sigma(k)},t_{\sigma(2n)})w_{\sigma(k)}
          %      \end{bmatrix}\\
          &=&  \begin{bmatrix}
                    \sum_{k=1}^n G(z_k,x_{\sigma(1)},t_k)w_k\\
                    \vdots\\
                    \sum_{k=1}^n G(z_k,x_{\sigma(n)},t_k)w_k
               \end{bmatrix}\\
          &=&  \begin{bmatrix}
                    F^{\sigma(1)}(z_1, \ldots, z_n,x_{1}, \ldots, x_{n})\\
                    \vdots\\
                    F^{\sigma(n)}(z_1, \ldots, z_n,x_{1}, \ldots, x_{n})
              \end{bmatrix}\\
          &=&   F(z_1, \ldots, z_n, x_1, \ldots, x_n) \cdot_2 \sigma. 
        \end{eqnarray*}
        This shows that $F$ is equivariant with respect to the actions. 
\end{proof} 

\end{document}
